\section{Related Work}\label{sec:rfw}

The idea of describing low-level data layout with high-level languages is
not new. The rich area of research makes it
challenging to fully contextualise our work within the
space. Simplifying our analysis a bit, we can roughly bifurcate the
literature into research on \emph{program synthesis} and \emph{program
abstraction}.

For instance, the PADS family of languages~\citep{FisherW11,
FisherG05, MandelbaumFWFG07}, PacketTypes~\citep{McCannC00},
Protege~\citep{WangG11}, DataScript~\citep{Back02}, Nail~\citep{Bangert_Zeldovich_14},
the generic
packet description by \citet{vanGeest_Swierstra_17}, the verified Protocol
Buffer~\citep{Ye_Delaware_19} built upon the \textsc{Narcissus}
framework~\citep{Delaware_SPYC_19}, EverParse~\citep{everparse},
and contiguity types~\citep{Slind_21} are all concerned
with synthesising a parser program (and also a pretty-printer for some of
them) from a high-level specification of the data format.

\dargent's primary focus is on the data refinement of algebraic
data types rather than securely operating on wire formats. Even though our
technology shares a lot in common, the problem we try to solve
is very different. In particular, \dargent is not a language for parsing or
converting between data formats. It is an extension to \cogent for its compiler to 
fine-tune the target code generation so that compiled code
is already in the desired format that is suitable for systems software.
In many cases, \dargent can eliminate the need for such a data marshalling tool entirely. 

Along with \dargent,
LoCal~\citep{Vollmer:2019}, SHAPES~\citep{Franco:17, Franco:19}
and \texttt{hobbit}~\citep{Diatchki_JL_05, Diatchki_Jones_06} fall in
 the program abstraction camp and 
are concerned with compiling
data structures in a program into specific layouts dictated by
the user. Programmers can therefore still work
with high-level source code, while the compiler 
does the heavy-lifting to generate the low-level mechanisms, retaining the 
separation of program logic from low-level concerns.

LoCal~\citep{Vollmer:2019} is a compiler for a first-order pure functional
language that can operate on recursive serialised data by translation into an
intermediate \emph{location} calculus, LoCal, mapping pointer indirections of
the high-level language to pointer arithmetic calculations on a base address.
The final compiler output is C code which, interestingly, preserves the
asymptotic complexity of the original recursive functions, although this
property is implementation-defined and not assured by any formal theorem.
Nonetheless, LoCal's type safety theorem does ensure a form of memory safety:
Each location is initialised and written to exactly once. The latter property
is a key difference to our work, since \Dargent can operate on mutable data by
virtue of \cogent's linear types. On the other hand, \cogent is a total
language and purposely lacks full support for recursion, we therefore do not
yet support recursive layout descriptions. Primitive recursive types for
\cogent are under development and use records~\citep{Murray:2019}. Since
we already support layout descriptions on records, we believe, once recursive
types are supported, adding support for recursive layouts would be a
straightforward engineering task. As a systems language, \cogent code often
uses abstract types such as arrays and  iteration constructs over such types.
Arrays and iteration constructs over arrays were recently verified through
\cogent's FFI~\cite{Cheung_OR_22}. We have ensured these proofs work with our
\dargent extensions. The most significant difference with LoCal is that
\Dargent is a \emph{certifying} data layout language, with generated theorems
that the translation is correct, whereas LoCal offers no verified guarantees
about its final compiler output.

The SHAPES~\citep{Franco:17, Franco:19} extension is designed for fine-tuning 
the layout of object classes for an object-oriented language to improve cache performance. 
It allows
users to define layout-unaware classes and specify what layout to use at
object instantiation time. This class parameterisation mechanism shares
some similarity with \dargent's layout polymorphism. The layouts that SHAPES
is concerned with are primarily arrays of values, which are key to better
cache locality but are not how compilers
of managed languages natively represent data in memory. Their layouts are not
down to the bit-level, but rather on the level of record fields.
In contrast, \dargent's layouts are lower-level and more flexible, and
are less tailored for a specific optimisation. SHAPES's type system maintains
memory safety properties of the program when it splits and lays out
boxed data types. Therefore it also plays a similar role as \cogent's
uniqueness type system. In our work, these two aspects are independently managed:
\dargent does not directly interfere with memory safety properties guaranteed
by \cogent's type system. 


The \texttt{hobbit} interpreter~\citep{Diatchki_JL_05} extends a Haskell-like functional language with
first-class support for bit-level types and operations (e.g.\ bit concatenation
and splitting), supporting external
representations of bit-level structures. Their work initially focused on
bit-data that can be stored within a single register and later gets extended
to memory areas realised as arrays~\citep{Diatchki_Jones_06}.
Instead of assigning a high-level type
and a low-level layout to an object in memory, their types already prescribe
the layouts, by virtue of the first-class bit-data support in the language. In
that sense, it is more comparable to the bit fields feature in the C language,
or to \cogent if the \dargent layout descriptors were subsumed by the \cogent types.
Their research novelty also lies in using advanced type system features that are readily
available in Haskell to encode the new language constructs and to perform
sophisticated typechecking, which is arguably an orthogonal
matter to data layouts.

Floorplan~\citep{Cronburg:2019} is also somewhat
relevant to \dargent, but hardly fits in either category. It is a
memory layout specification language for declaratively describing the
structure of a heap as laid out by a memory manager.
It therefore chiefly serves the implementors of memory managers rather than
systems developers and users in general, and the abstraction it provides
does not necessarily extends to algebraic types of heap objects.
The compiler follows the specification to generate memory-safe Rust code to perform
common tasks that are needed in the implementation of a memory manager.
The semantics of a heap layout
specification is denoted by the set of values that the heap can take.
In contrast, the semantics of a \dargent layout is characterised by the getter
and the setter functions.

